\section{The observations needed to select a winner}
\label{sec:certificate}
View the votes as a table with one row per candidate and one column per restoration. A query reveals one count. In the following result, unread counts may independently take any integer from zero to the number of panels \(B\). This is a model of the information available to the search. A fitted classifier can impose further constraints on its predictions.

\begin{proposition}
After any queries, let \(u_i\) be the smallest count read in row \(i\), initially \(B\). Set \(\ell_i=u_i\) for a fully read row and \(\ell_i=0\) otherwise. The observations determine \(w\) as the first maximum precisely when
\[
  u_i<\ell_w\quad(i<w),\qquad u_i\le\ell_w\quad(i>w).
\]
With a fixed restoration order within each row, best-first nomination uses the fewest queries that can establish the winner on the realized table whenever \(w>1\) or some rival has a positive final score.
\end{proposition}

\begin{proof}
Every compatible row minimum lies between \(\ell_i\) and \(u_i\). Both endpoints can be attained independently across rows: fill unread entries with zero for the lower endpoint and \(B\) for the upper. Putting the winner at its lower endpoint and every rival at its upper endpoint gives the two inequalities, including the strict one for earlier rivals.

Write \(s=v(S_w)\). Under the stated condition any sufficient collection of prefixes must exhaust row \(w\): with an unread entry its lower endpoint is zero, which cannot defeat an earlier rival strictly or a positive rival weakly. Each earlier rival must be read until its upper bound first falls below \(s\), and each later rival until its bound first reaches or falls below \(s\). A prefix of length zero can suffice. The true winner's running bound always remains at least \(s\). Once a rival reaches its required cutoff it ranks behind the winner and the search never reads it again. On termination the inequalities hold, so the search reads exactly the required prefixes.
\end{proof}

The exception is observable in hindsight. With two rows of length \(H\), let the first contain only \(B\) and the second only zeros. Best-first reads all \(H\) entries of the first row. Reading just the first zero in the second row also proves the first candidate wins. If both rows instead contain only \(B\), the unique shortest proof reads the first row alone. For \(H>1\) no first query starts a shortest proof on both tables. These counts concern votes read in fixed order. A cached reference prediction can serve several votes, and computation time also depends on the learner.
